\chapter{Mennesket som feste}

\begin{motto}
Smak skifter med moten.\\ Korleis ei hand finn ein dørvrider, gjer det ikkje.
\end{motto}

Av alle skuldingane mot designfaget er den om subjektivitet den mest hardnakka. Design er smak, heiter det; det handlar om kva som ser fint ut, og om kva som sel, og smak har inga sanning å svare for. Dersom det var heile historia, ville faget verkeleg stå utan feste. Men det er ikkje heile historia, og dette kapitlet handlar om kvifor. Ein gong i andre helvta av det tjuande hundreåret skjedde det noko med designfaget som flytta det vekk frå smaken og mot noko fastare: brukarsentreringa gav design eit objektivt feste, nemleg mennesket. Affordans, signifier, kartlegging, tilbakemelding — dette er ikkje synsmåtar om kva som er pent. Det er overførbar, kumulativ kunnskap om korleis menneske møter ting, og den kunnskapen lèt seg prøve, korrigere og byggje vidare på. Mennesket sit ikkje fast i moten. Difor kan eit fag som tek mennesket som mål, byggje på fast grunn.

\section{Frå smak til menneske}

Vendinga har eit ansikt og ein tittel. Då Donald Norman gav ut \emph{The Psychology of Everyday Things} i 1988 — seinare omdøypt til \emph{The Design of Everyday Things} \autocite{norman2013} — gjorde han noko meir enn å skrive endå ei bok om god form. Han flytta grunnlaget for designdommen. Spørsmålet var ikkje lenger «likar eg dette?», men «greier eit menneske å bruke dette utan å bli forvirra, og kvifor eller kvifor ikkje?». Det er eit anna slag spørsmål, fordi det har eit svar som ikkje er mitt eller ditt, men som ligg i møtet mellom ein kropp, ein hjerne og ein ting. Den umoglege døra som du dreg i når du skal dytte, er ikkje stygg. Ho er feildesigna, og feilen lèt seg påvise utan å spørje om smak: ho gjev feil signal om kva ein skal gjere.

Norman henta omgrepa sine frå psykologien og frå persepsjonsforskaren James Gibson, og det er nett opphavet som gjev dei feste. \emph{Affordans} — Gibsons ord for det handlingsrommet ein ting tilbyr ein organisme — vart hos Norman eit verktøy for å snakke presist om kvifor ein flate inviterer til å bli dytta og ein knott til å bli vridd. Affordans er ikkje noko designaren finn opp; det er noko som finst i forholdet mellom tingen og den som brukar han, og som difor kan studerast utanfrå. Ein trapp gjev affordans for å gå opp same kva eg synest om henne. Ei sakshandsamingsside som skjuler den eine knappen brukaren leitar etter, sviktar uavhengig av om fargane er smakfulle. Med affordansomgrepet fekk faget eit grep som peikar bort frå designaren si eiga oppleving og mot ein eigenskap ved sjølve situasjonen.

At opphavet er psykologien, er ikkje ein tilfeldig biografisk detalj; det er sjølve grunnen til at omgrepa held. Då Norman bytte tittelen frå \emph{psykologien} til \emph{design}, signaliserte han ikkje at faget hadde gjeve opp vitskapen, men at innsiktene frå persepsjon, minne og merksemd no var modne nok til å bere ein praksis. Eit dårleg utforma kontrollpanel er ikkje dårleg fordi nokon mislikar det; det er dårleg fordi det krev meir av korttidsminnet enn eit menneske har, eller fordi det bryt med dei mønstera auget ventar. Slike påstandar svarar til noko utanfor designaren — til kor mange einingar eit menneske kan halde i hovudet på éin gong, til kor fort blikket finn ein kontrast. Det er kunnskap av same slag som ergonomien sin: ho gjeld for kroppen og sansane som dei faktisk er, ikkje for kva ein generasjon synest er moderne. Difor lèt ho seg òg overføre over landegrenser og over tiår. Ein dørvrider som er lett å gripe i Tokyo, er lett å gripe i Trondheim, fordi handa er den same.

\section{Eit lite, hardt vokabular}

Det Norman bygde, var ikkje ein enkelt innsikt, men eit lite vokabular av samverkande omgrep, og styrken ligg i at dei heng saman til ein samanhengande modell av interaksjon. Ved sida av affordansen står \emph{kartlegginga}: tilhøvet mellom kontrollar og verknader, mellom brytaren og lampa, mellom rattet og hjula. Ei god kartlegging er den der det romlege og det logiske mønsteret i kontrollane svarar til mønsteret i det dei styrer, slik at ein kokeplate-panelet «seier» kva knapp som høyrer til kva plate. Dette er ikkje pynt; det er ein påvisbar skilnad mellom oppsett som folk meistrar med ein gong og oppsett som tvingar dei til å lære utanåt eller prøve seg fram. Ved sida av kartlegginga står \emph{tilbakemeldinga}: at systemet seier frå om kva det har gjort, slik at brukaren veit om handlinga gjekk gjennom. Eit grensesnitt utan tilbakemelding let mennesket stå att i uvisse, og uvissa er ikkje ei smakssak.

Det avgjerande for argumentet i denne boka er ein detalj i revisjonshistoria. I 2013-utgåva la Norman til eit omgrep han ikkje hadde hatt før: \emph{signifier} \autocite{norman2013}. Affordansar finst, sa han no, men dei hjelper berre om dei er \emph{synlege} — og det synlege teiknet som fortel brukaren kva han kan gjere og kvar, det er signifieren. Skiljet kan verke som ei finurlegheit, men det er det motsette av tomt: det rettar opp ein systematisk mistydnad av den fyrste utgåva, der lesarar hadde teke til å kalle kvart vesle visuelle hint ein «affordans». At ein forfattar går attende til sitt eige sentrale omgrep, ser at det vart misforstått, og finskil det i to — affordansen som høyrer til forholdet, signifieren som høyrer til kommunikasjonen — er ikkje teikn på at faget flyt. Det er teikn på det motsette. Det er slik kunnskap går framover: ved at omgrepa blir skarpare med bruk, at feil bruk blir oppdaga, og at den neste utgåva veit meir enn den førre. Eit felt utan fundament rettar ikkje vokabularet sitt. Det berre skiftar mote.

Norman følgde sjølv opp med å vise at det funksjonelle og det kjenslemessige ikkje står mot kvarandre. I \emph{Emotional Design} \autocite{norman2004emotional} la han fram at produkt verkar på fleire plan på éin gong — det umiddelbart sanselege, det åtferdsmessige, det reflekterande — og at vakre ting faktisk fungerer betre, fordi velvære gjer menneske meir tolsame og meir kreative i møtet med små feil. Dette er ikkje ei innrømming av at design likevel er smak. Det er ei utviding av det same prosjektet: å gjere greie for menneskets faktiske møte med ting, no også på det affektive planet, som ein eigenskap som lèt seg studere snarare enn berre kjennast.

\section{Det målbare og det konkrete}

Om Norman gav faget eit vokabular, gav andre det måleskeier og framgangsmåtar. Jakob Nielsen sitt arbeid med \emph{usability engineering} \autocite{nielsen1993} er det tydelegaste dømet på at brukarsentreringa ikkje stoppa ved kloke ord, men gjekk vidare til noko ein kan telje. Brukskvalitet, hos Nielsen, er ikkje ein eigenskap ein anar; det er ei mengd storleikar ein måler: kor lang tid ein oppgåve tek, kor mange feil brukaren gjer, kor mykje som hugsast ved neste gong, kor mange som lukkast i det heile. At desse storleikane lèt seg måle, gjer at to fagfolk kan vere usamde om eit grensesnitt og så \emph{avgjere} usemja ved å teste det på menneske. Det er nett den evna kritikaren av designfaget hevdar at faget manglar — evna til å ta feil og bli korrigert av verda. På feltet for brukskvalitet finst denne evna. Ein påstand om at den nye utforminga er betre, kan visast å vere gal.

Nielsen sine ti heuristikkar for brukargrensesnitt er eit anna døme på kumulasjon i praksis: destillerte, overførbare retningslinjer som ein ny utøvar kan lære og bruke på eit grensesnitt forfattaren aldri såg. Dei er ikkje lovar i fysikkens forstand, og dei gjev ikkje eitt rett svar. Men dei er heller ikkje smak. Dei er komprimert røynsle om kvar menneske snublar, av same slaget som ein lege sine kliniske teikn eller ein ingeniør sine tommelfingerreglar — kunnskap som er prøvd mot mange tilfelle og funnen å halde i det breie. At eit felt produserer slike retningslinjer, og at dei blir reviderte når nye studiar viser kvar dei sviktar, er ikkje fråvêret av eit fundament. Det er fundamentet i arbeid.

Det er verdt å dvele ved kor sterk denne forma for prøving eigentleg er, for ho råkar kritikken på det ømmaste punktet. \textsc{Grunnlaust} sin hardaste påstand er at design ikkje kan ta feil — at faget manglar den evna til å bli motsagt av verda som skil ein vitskap frå ei meining. Men ein brukstest \emph{er} ei prøving som kan gå imot den som utforma. Designaren går inn i rommet overtydd om at det nye oppsettet er klarare; fem menneske set seg ned, og tre av dei finn ikkje fram. Då er saka avgjord, ikkje av smak og ikkje av makt, men av at verda svarte nei. Nielsen synte til og med at talet på testpersonar som trengst for å avdekkje dei fleste feila er lite og lèt seg rekne ut — eit funn om sjølve metoden, ikkje berre om eit grensesnitt. Eit felt som forskar på sine eigne framgangsmåtar og finn ut kor mange forsøk som skal til før ein veit nok, driv ikkje med smak. Det driv med metode, og metoden har den eigenskapen kritikaren hevda at faget mangla: han kan vise at den profesjonelle tok feil.

Alan Cooper gav vendinga endå ei form med \emph{personas} \autocite{cooper1999inmates}. Her er ein metode for å halde brukaren føre seg gjennom heile utviklinga: ein presis, namngjeven, oppdikta, men forskingsfunder modell av den ein utformar for, slik at avgjerder kan prøvast mot eit konkret menneske i staden for mot ein tåkete «alle». Cooper sitt poeng var skarpt: når ingen er ansvarleg for å representere brukaren, vinn maskina og ingeniøren sine vanar, og resultatet blir produkt som «driv oss til vanvit». Personas er svaret — ein disiplin for å gjere det abstrakte mennesket konkret nok til å kunne seie nei på dets vegne. Her må vi vere ærlege om ei reell fare, og kapitlet om misbruk kjem attende til henne: ein persona kan òg bli ein konstruksjon i teneste for oppdragsgjevaren, ein figur dikta opp for å rettferdiggjere det leiinga alt hadde bestemt. Det hender, og det er gale når det hender. Men ein metode som lèt seg misbruke, er ikkje difor utan verd; skalpellen blir ikkje verdlaus av at nokon kan skade med han. Misbruket er ein kritikk av utøvaren, ikkje av grunnlaget. At brukaren kan forfalskast, føreset nett at det finst ein ekte brukar å forfalske mot — og det er den ekte brukaren faget tek som sitt feste.

\section{Då brukaren vart eit seriøst objekt}

Bak verktøya ligg ei forsking som gjorde mennesket til eit objekt verdig vitskapleg alvor, og det er denne forskinga som hindrar at brukarsentreringa kan avfeiast som folkeleg psykologi. Lucy Suchman sin studie \emph{Plans and Situated Actions} \autocite{suchman1987} er eit vendepunkt. Ho studerte korleis folk faktisk strevde med ein avansert kopimaskin, og synte at menneskeleg handling ikkje er gjennomføringa av ein ferdig plan, slik den tidlege kognitivismen og kunstig-intelligens-forskinga hadde gått ut frå, men noko \emph{situert} — laga i augneblinken, i respons på ein konkret og rotete situasjon. Det er ei empirisk innsikt med harde følgjer for design: eit system som føreset at brukaren følgjer ein plan, vil svikte mennesket som faktisk improviserer. Suchman gjorde brukaren til gjenstand for streng, etnografisk informert gransking, og endra med det kva ein måtte vite for å utforme noko menneske skal bruke.

Same rørsla finn vi hos Terry Winograd og Fernando Flores, som i \emph{Understanding Computers and Cognition} \autocite{winograd1986} henta fenomenologi og språkfilosofi inn i datavitskapen og argumenterte for at design måtte ta utgangspunkt i korleis menneske faktisk er i verda — i praksis, i språk, i samhandling — og ikkje i ein abstrakt modell av informasjonshandsaming. Boka sin undertittel, eit nytt \emph{fundament} for design, var ikkje tilfeldig vald. Det interaksjonen mellom menneske og maskin trong, hevda dei, var ikkje meir reknekraft, men eit sannare bilete av kva mennesket er — eit vesen som alltid alt er i ein samanheng, med eit språk og ein praksis, før det i det heile møter ein skjerm. Og Paul Dourish samla denne tråden i \emph{Where the Action Is} \autocite{dourish2001}, der han bygde ut omgrepet om \emph{embodied interaction}: at meininga i eit grensesnitt blir til gjennom kroppsleg og sosial handling, ikkje berre i hovudet. Vi forstår ein ting ved å handle med han, og handlinga er kroppsleg og situert; difor kan ikkje eit grensesnitt utformast godt utan kunnskap om korleis kroppen og det sosiale faktisk verkar. Til saman utgjer Suchman, Winograd og Flores, og Dourish ein forskingstradisjon — med teoriar, metodar, usemjer og framsteg — om mennesket i møte med det utforma. Det er ikkje smak. Det er kunnskap, og det er kumulativ: kvar bygde på, og kvar korrigerte, dei føre. Suchman gav grunnlaget om det situerte; Winograd og Flores gav det filosofiske fundamentet; Dourish samla det til ein teori om den kroppslege interaksjonen. Slik veks ein kunnskap, ledd for ledd, og resultatet er ein arv ein ny utøvar kan ta i mot, ikkje ein smak ho må gisse seg til.

Sett dette saman, og påstanden om at design berre er smak, fell. Brukarsentreringa gav faget eit objektivt feste fordi mennesket er felles. Ei hand som ikkje finn knappen, ein hugs som sviktar under press, eit blikk som søkjer signifieren og ikkje finn han — dette er ikkje meiningar. Det er tilhøve i verda, og dei lèt seg studere, måle og betre. At det òg finst smak i faget, og at smak rår der mennesket sine grenser ikkje gjev svar, er sant. Men under smaken ligg eit lag av overførbar, korrigerbar kunnskap om menneske og ting, og det laget er fast. Eit fag som kviler på det, kviler ikkje på sand.

\vignett

Men mennesket er ikkje berre ein brukar med hender og blikk. Det er òg ein skapnad som lever i ei verd faget sjølv er med på å lage. Det fører oss til det djupaste forsvaret av alle.

\vidarelesing{\fullcite{norman2013}; \fullcite{nielsen1993}; \fullcite{cooper1999inmates}; \fullcite{suchman1987}; \fullcite{winograd1986}; \fullcite{dourish2001}; \fullcite{norman2004emotional}.}
